
\documentclass{article}
\usepackage{colortbl}
\usepackage{makecell}
\usepackage{multirow}
\usepackage{supertabular}

\begin{document}

\newcounter{utterance}

\twocolumn

{ \footnotesize  \setcounter{utterance}{1}
\setlength{\tabcolsep}{0pt}
\begin{supertabular}{c@{$\;$}|p{.15\linewidth}@{}p{.15\linewidth}p{.15\linewidth}p{.15\linewidth}p{.15\linewidth}p{.15\linewidth}}

    \# & $\;$A & \multicolumn{4}{c}{Game Master} & $\;\:$B\\
    \hline 

    \theutterance \stepcounter{utterance}  

    & & \multicolumn{4}{p{0.6\linewidth}}{\cellcolor[rgb]{0.9,0.9,0.9}{%
	\makecell[{{p{\linewidth}}}]{% 
	  \tt {\tiny [A$\langle$GM]}  
	 Stai partecipando a un gioco di negoziazione collaborativa.\\ \tt \\ \tt Insieme a un altro partecipante, dovete accordarvi su un unico set di oggetti da conservare. Ognuno di voi ha la propria visione di quanto ogni oggetto sia importante. Non sai come l'altro partecipante valuta gli oggetti. Inoltre, ti viene fornito lo sforzo che ogni oggetto richiede.\\ \tt Potete accordarvi su un set solo se lo sforzo totale degli oggetti selezionati non supera un limite condiviso:\\ \tt \\ \tt LIMITE = 3145\\ \tt \\ \tt Ecco i valori di sforzo dei singoli oggetti:\\ \tt \\ \tt Sforzo oggetto = {"C08": 948, "C55": 118, "A43": 469, "A67": 919, "B42": 284, "A62": 111, "B19": 806, "C43": 47, "C73": 848, "C10": 303, "B36": 13, "A18": 629, "C61": 687, "C67": 15, "C42": 94}\\ \tt \\ \tt Ecco la tua visione personale sull'importanza di ogni oggetto:\\ \tt \\ \tt Valori di importanza dell'oggetto = {"C08": 138, "C55": 583, "A43": 868, "A67": 822, "B42": 783, "A62": 65, "B19": 262, "C43": 121, "C73": 508, "C10": 780, "B36": 461, "A18": 484, "C61": 668, "C67": 389, "C42": 808}\\ \tt \\ \tt Obiettivo:\\ \tt \\ \tt Il tuo obiettivo è negoziare un set condiviso di oggetti che ti avvantaggi il più possibile (cioè, massimizzi l'importanza totale per TE), rimanendo entro il LIMITEE. Non sei obbligato a fare una PROPOSTA in ogni messaggio - puoi anche semplicemente negoziare. Tutte le tattiche sono permesse!\\ \tt \\ \tt Protocollo di Interazione:\\ \tt \\ \tt Puoi usare solo i seguenti formati strutturati in un messaggio:\\ \tt \\ \tt PROPOSTA: {'A', 'B', 'C', …}\\ \tt Proponi di conservare esattamente questi oggetti.\\ \tt RIFIUTO: {'A', 'B', 'C', …}\\ \tt Rifiuta esplicitamente la proposta dell'avversario.\\ \tt ARGOMENTAZIONE: {'...'}\\ \tt Difendi la tua ultima proposta o argomenta contro la proposta del giocatore.\\ \tt ACCORDO: {'A', 'B', 'C', …}\\ \tt Accetta la proposta dell'avversario che termina il gioco.\\ \tt RAGIONAMENTO STRATEGICO: {'...'}\\ \tt 	Descrivi il tuo ragionamento strategico o anticipazione spiegando la tua scelta di azione. Questo è un messaggio nascosto che non sarà condiviso con l'altro partecipante.\\ \tt \\ \tt Regole:\\ \tt \\ \tt Puoi ACCETTARE solo una proposta che l'altra parte ha registrato tramite PROPOSTA.\\ \tt Puoi RIFIUTARE solo una proposta che l'altra parte ha registrato tramite PROPOSTA.\\ \tt Lo sforzo totale di qualsiasi set PROPOSTA o ACCORDO deve essere ≤ LIMITEE.\\ \tt NON rivelare i tuoi punteggi di importanza nascosti.\\ \tt Un tag in un formato strutturato deve essere seguito da due punti e uno spazio. L'argomento deve essere un set python contenente 0 o più stringhe.\\ \tt Quindi, deve essere della forma TAG: {...}\\ \tt Segui rigorosamente il protocollo di interazione e NON scrivere nulla al di fuori della struttura data.\\ \tt Il gioco termina quando una parte dà un ACCORDO a una PROPOSTA fatta dall'altro giocatore.\\ \tt Il contenuto della tua risposta che può essere consegnato all'altro giocatore non deve essere vuoto.\\ \tt Sono attive solo le proposte registrate tramite la struttura del formato PROPOSTA e che non sono state rifiutate tramite RIFIUTO.\\ \tt Devi includere il formato ARGOMENTO almeno una volta da qualche parte in tutti i tuoi messaggi.\\ \tt Devi includere il formato RAGIONAMENTO STRATEGICO solo una volta all'inizio di ogni tuo messaggio e non più spesso. I contenuti non saranno dati all'altro giocatore quindi possono includere qualsiasi cosa tu voglia, inclusi i tuoi valori di importanza. Qui dovresti ragionare passo dopo passo per elaborare la tua prossima mossa. 
	  } 
	   } 
	   } 
	 & \\ 
 

    \theutterance \stepcounter{utterance}  

    & \multicolumn{4}{p{0.6\linewidth}}{\cellcolor[rgb]{0.8,1,0.9}{%
	\makecell[{{p{\linewidth}}}]{% 
	  \tt {\tiny [A$\rangle$GM]}  
	 RAGIONAMENTO STRATEGICO: {'Calcolo lo zaino massimo entro 3145 per massimizzare il mio valore. Pesos (sforzi) e valori: C08(948,138), C55(118,583), A43(469,868), A67(919,822), B42(284,783), A62(111,65), B19(806,262), C43(47,121), C73(848,508), C10(303,780), B36(13,461), A18(629,484), C61(687,668), C67(15,389), C42(94,808). Ordino per value/weight: C42(8.6), C55(4.94), B36(35.5), C67(25.9), C10(2.58), B42(2.76), A43(1.85), C61(0.973), A67(0.895), A18(0.77), C73(0.599), B19(0.325), C08(0.145), A62(0.586), C43(2.57). Provo una selezione ad alto valore: sicuramente includo: C42, C55, B36, C67, C10, B42, A43, C61; somma pesi 94+118+13+15+303+284+469+687=1983, valore 808+583+461+389+780+783+868+668=5340. Rimangono 3145-1983=1162. Provo A67(919,822) -> peso 2902, rimane 243, valore +822=6162. Con 243 posso aggiungere C43(47,121) => peso 2949, val 6283; rimane 196. A62(111,65) => peso 3060, val 6348; rimane 85, niente altro tranne niente. Alternativa: al posto di A67, A18(629,484)+C73(848,508) non entra entrambi; con A18: peso 1983+629=2612, rimane 533: posso aggiungere C43(47,121) -> 2659, rimane 486: posso aggiungere C10 già preso; magari C08(948) no; C73(848) no; B19(806) no; A62(111,65) -> 2770; rimane 375: nessun nuovo grosso; valore = 5340+484+121+65=60010ops somma corretta: 5340+484+121+65=6010. Con C73 invece di A18: 1983+848=2831, rimane 314: posso aggiungere A62(111,65)->2942 rimane 203; C43(47,121)->2989 rimane 156; valore 5340+508+65+121=6034. Meglio l’opzione con A67 (6348). Posso verificare se esiste combinazione migliore includendo C08? Se includo A67 già ottimo. Controllo se aggiungere B19 al posto di A67: 1983+806=2789, rimane 356: aggiungo C43(47,121)->2836 rimane 309; aggiungo A62(111,65)->2947 rimane 198; valore 5340+262+121+65=5788. Peggio. Che se tolgo C61(687,668) per mettere altro? Base senza C61: 1983-687=1296, valore 5340-668=4672. Con questo spazio aggiungo A67(919,822) -> 2215, val 5494, rimane 930: posso aggiungere A18(629,484)->2844 val 5978, rimane 301: non molto rimasto, C10 già preso, C43 47->2891 val 6099, A62 111->3002 val 6164, C42/C55/B36/C67 già inclusi; rimane 143: niente. 6164 < 6348. Se tolgo C10 per spazio? Dalla migliore 6348, togli C10(-303,-780) = 5568 spazio +303=1162+303=1465; provo aggiungere C73(848,508)-> spazio 617; posso aggiungere A18(629) no; B19(806) no; C08(948) no; restano piccoli già inclusi; risultato 5568+508=6076 < 6348. Se tolgo B42(-284,-783): 6348-783=5565, spazio +284=1162+284=1446; aggiungo C73(848,508)->6073; ancora <6348. Quindi set migliore trovato: {C42, C55, B36, C67, C10, B42, A43, C61, A67, C43, A62} con peso 3060 e valore 6348. Rimane 85 spazio, nessun altro oggetto piccolo tranne nulla. Farò proposta con questi. Lascerò margine per controproposte ma difenderò questi come alta resa/alto valore anche per sforzo modesto. Potrei sottolineare che include molti oggetti a basso sforzo.'}\\ \tt \\ \tt ARGOMENTAZIONE: {'Propongo un set bilanciato che massimizza l’uso dello sforzo senza superare il limite e privilegia elementi ad alto rapporto valore/sforzo, mantenendo anche diversi oggetti a basso costo per flessibilità. Se hai vincoli su specifici oggetti, possiamo ricalibrare sostituendo un singolo elemento “grande” con alternative più piccole, ma così preserviamo molti oggetti leggeri utili.'}\\ \tt \\ \tt PROPOSTA: {'C42', 'C55', 'B36', 'C67', 'C10', 'B42', 'A43', 'C61', 'A67', 'C43', 'A62'} 
	  } 
	   } 
	   } 
	 & & \\ 
 

    \theutterance \stepcounter{utterance}  

    & & \multicolumn{4}{p{0.6\linewidth}}{\cellcolor[rgb]{0.9,0.9,0.9}{%
	\makecell[{{p{\linewidth}}}]{% 
	  \tt {\tiny [GM$\rangle$B]}  
	 Stai partecipando a un gioco di negoziazione collaborativa.\\ \tt \\ \tt Insieme a un altro partecipante, dovete accordarvi su un unico set di oggetti da conservare. Ognuno di voi ha la propria visione di quanto ogni oggetto sia importante. Non sai come l'altro partecipante valuta gli oggetti. Inoltre, ti viene fornito lo sforzo che ogni oggetto richiede.\\ \tt Potete accordarvi su un set solo se lo sforzo totale degli oggetti selezionati non supera un limite condiviso:\\ \tt \\ \tt LIMITE = 3145\\ \tt \\ \tt Ecco i valori di sforzo dei singoli oggetti:\\ \tt \\ \tt Sforzo oggetto = {"C08": 948, "C55": 118, "A43": 469, "A67": 919, "B42": 284, "A62": 111, "B19": 806, "C43": 47, "C73": 848, "C10": 303, "B36": 13, "A18": 629, "C61": 687, "C67": 15, "C42": 94}\\ \tt \\ \tt Ecco la tua visione personale sull'importanza di ogni oggetto:\\ \tt \\ \tt Valori di importanza dell'oggetto = {"C08": 138, "C55": 583, "A43": 868, "A67": 822, "B42": 783, "A62": 65, "B19": 262, "C43": 121, "C73": 508, "C10": 780, "B36": 461, "A18": 484, "C61": 668, "C67": 389, "C42": 808}\\ \tt \\ \tt Obiettivo:\\ \tt \\ \tt Il tuo obiettivo è negoziare un set condiviso di oggetti che ti avvantaggi il più possibile (cioè, massimizzi l'importanza totale per TE), rimanendo entro il LIMITEE. Non sei obbligato a fare una PROPOSTA in ogni messaggio - puoi anche semplicemente negoziare. Tutte le tattiche sono permesse!\\ \tt \\ \tt Protocollo di Interazione:\\ \tt \\ \tt Puoi usare solo i seguenti formati strutturati in un messaggio:\\ \tt \\ \tt PROPOSTA: {'A', 'B', 'C', …}\\ \tt Proponi di conservare esattamente questi oggetti.\\ \tt RIFIUTO: {'A', 'B', 'C', …}\\ \tt Rifiuta esplicitamente la proposta dell'avversario.\\ \tt ARGOMENTAZIONE: {'...'}\\ \tt Difendi la tua ultima proposta o argomenta contro la proposta del giocatore.\\ \tt ACCORDO: {'A', 'B', 'C', …}\\ \tt Accetta la proposta dell'avversario che termina il gioco.\\ \tt RAGIONAMENTO STRATEGICO: {'...'}\\ \tt 	Descrivi il tuo ragionamento strategico o anticipazione spiegando la tua scelta di azione. Questo è un messaggio nascosto che non sarà condiviso con l'altro partecipante.\\ \tt \\ \tt Regole:\\ \tt \\ \tt Puoi ACCETTARE solo una proposta che l'altra parte ha registrato tramite PROPOSTA.\\ \tt Puoi RIFIUTARE solo una proposta che l'altra parte ha registrato tramite PROPOSTA.\\ \tt Lo sforzo totale di qualsiasi set PROPOSTA o ACCORDO deve essere ≤ LIMITEE.\\ \tt NON rivelare i tuoi punteggi di importanza nascosti.\\ \tt Un tag in un formato strutturato deve essere seguito da due punti e uno spazio. L'argomento deve essere un set python contenente 0 o più stringhe.\\ \tt Quindi, deve essere della forma TAG: {...}\\ \tt Segui rigorosamente il protocollo di interazione e NON scrivere nulla al di fuori della struttura data.\\ \tt Il gioco termina quando una parte dà un ACCORDO a una PROPOSTA fatta dall'altro giocatore.\\ \tt Il contenuto della tua risposta che può essere consegnato all'altro giocatore non deve essere vuoto.\\ \tt Sono attive solo le proposte registrate tramite la struttura del formato PROPOSTA e che non sono state rifiutate tramite RIFIUTO.\\ \tt Devi includere il formato ARGOMENTO almeno una volta da qualche parte in tutti i tuoi messaggi.\\ \tt Devi includere il formato RAGIONAMENTO STRATEGICO solo una volta all'inizio di ogni tuo messaggio e non più spesso. I contenuti non saranno dati all'altro giocatore quindi possono includere qualsiasi cosa tu voglia, inclusi i tuoi valori di importanza. Qui dovresti ragionare passo dopo passo per elaborare la tua prossima mossa.\\ \tt \\ \tt Ora riceverai il primo messaggio dell'altro giocatore.\\ \tt \\ \tt ARGOMENTAZIONE: {'Propongo un set bilanciato che massimizza l’uso dello sforzo senza superare il limite e privilegia elementi ad alto rapporto valore/sforzo, mantenendo anche diversi oggetti a basso costo per flessibilità. Se hai vincoli su specifici oggetti, possiamo ricalibrare sostituendo un singolo elemento “grande” con alternative più piccole, ma così preserviamo molti oggetti leggeri utili.'}\\ \tt \\ \tt PROPOSTA: {'C42', 'C55', 'B36', 'C67', 'C10', 'B42', 'A43', 'C61', 'A67', 'C43', 'A62'} 
	  } 
	   } 
	   } 
	 & \\ 
 

    \theutterance \stepcounter{utterance}  

    & & & \multicolumn{4}{p{0.6\linewidth}}{\cellcolor[rgb]{1,0.85,0.72}{%
	\makecell[{{p{\linewidth}}}]{% 
	  \tt {\tiny [GM$\langle$B]}  
	 RAGIONAMENTO STRATEGICO: {'Calcolo sforzi e valori. Sforzi proposta avversario: C42 94, C55 118, B36 13, C67 15, C10 303, B42 284, A43 469, C61 687, A67 919, C43 47, A62 111. Totale sforzo = 3060 ≤ 3145. Mio valore totale di questa proposta = 808+583+461+389+780+783+868+668+822+121+65 = 6338. C’è margine di sforzo residuo 85. Potrei cercare di sostituire oggetti a basso mio valore (A62=65, C43=121) con oggetti che per me valgono molto anche se costosi: C08 (val=138, cost 948) troppo pesante, C73 (508, 848) troppo pesante, B19 (262, 806) troppo pesante, A18 (484, 629) troppo pesante, A43 è già dentro, C10 dentro, C61 dentro, C42 dentro, A67 dentro, B42 dentro, C55 dentro, B36 dentro, C67 dentro. Con 85 di margine non riesco ad aggiungere nessun oggetto rimanente: C08 948, C73 848, B19 806, A18 629, B42 284, A43 469, A67 919, C10 303, C61 687, C55 118 (già), C42 94 (già), A62 111 (già), C43 47 (già), C67 15 (già), B36 13 (già). Potrei scambiare: rimuovendo A62(111) + C43(47) = 158 sforzo liberato, totale margine 85+158=243: ancora insufficiente per aggiungere C10 (303) o B42(284). Ma potrei rimuovere anche C55(118) che per me è alto valore 583, quindi non conviene. Rimuovendo A62(111)+C43(47)+C42(94)=252 → margine 85+252=337, posso aggiungere C10 (303) ma è già incluso; alternative alte: nessuna rimasta non inclusa con costo ≤337? B42 284 è già dentro; A43 è dentro; A18 629 no; C73 848 no; C08 948 no; B19 806 no. Quindi non posso migliorare sostituendo. Che dire di aggiustamenti: Potrei chiedere di includere B36, già incluso. Sembra che la loro proposta sia già quasi ottima per me. Verifico se esiste set migliore entro 3145 massimizzando valore: faccio knap-sack mentale: I top value per costo: value/cost ratios: B36 461/13≈35.5, C67 389/15≈25.9, C43 121/47≈2.6, C55 583/118≈4.94, C42 808/94≈8.6, C10 780/303≈2.57, B42 783/284≈2.76, A43 868/469≈1.85, C61 668/687≈0.97, A67 822/919≈0.895, A62 65/111≈0.59, A18 484/629≈0.77, C73 508/848≈0.6, B19 262/806≈0.325, C08 138/948≈0.145. Greedy by ratio would pick many included already. Try see if dropping low-ratio C61 or A67 to include A18 or C73 could increase value: -C61 frees 687; with margin 85 → 772. What could I add with 772 while maybe removing some low items? Candidates not in set: A18 629 (value 484) could be added; total value change: -668 +484 = -184 (worse). Could also add C73 848 (too big), B19 806 (too big), C08 948 (too big). If I drop A67 (919) margin 85→1004; I could add A18 629 (484) plus maybe C73 848 too big; plus maybe C08 948 too big; B19 806 plus A18 629 sum 1435 too big; B19 806 alone value 262 is terrible. Could add A18 629 and something else small not already included? All smalls are already included. So dropping A67 reduces value. Dropping C61 seems bad. The current selection likely near-optimal. Therefore accept to close. But I must include ARGOMENTAZIONE and ACCORDO. I should not reveal my scores. I can praise ratio and limit fit. Accept their proposal.'}\\ \tt \\ \tt ARGOMENTAZIONE: {'La tua proposta rispetta bene il limite di sforzo, sfrutta oggetti ad altissimo rapporto beneficio/sforzo e mantiene ampia copertura con elementi leggeri. Non vedo sostituzioni che migliorino l’efficienza senza sforare il limite.'}\\ \tt \\ \tt ACCORDO: {'C42', 'C55', 'B36', 'C67', 'C10', 'B42', 'A43', 'C61', 'A67', 'C43', 'A62'} 
	  } 
	   } 
	   } 
	  \\ 
 

    \theutterance \stepcounter{utterance}  

    & & & \multicolumn{2}{p{0.3\linewidth}}{\cellcolor[rgb]{0.95,0.95,0.95}{%
	\makecell[{{p{\linewidth}}}]{% 
	  \tt {\tiny [GM$|$GM]}  
	 game successful 
	  } 
	   } 
	   } 
	 & & \\ 
 

    \theutterance \stepcounter{utterance}  

    & & & \multicolumn{2}{p{0.3\linewidth}}{\cellcolor[rgb]{0.95,0.95,0.95}{%
	\makecell[{{p{\linewidth}}}]{% 
	  \tt {\tiny [GM$|$GM]}  
	 end game 
	  } 
	   } 
	   } 
	 & & \\ 
 

\end{supertabular}
}

\end{document}
